\section{Connecting orientations by flips}\label{sec:connectivity}

In this section we prove Theorem~\ref{thm:connectivity}.
We first move the edges of an orientation into a normal form determined by the
corners, Section~\ref{subs:preparation}.
Section~\ref{subs:slides} constructs explicit paths of flips which exchange two
edges and restore all others.
With these we treat complete graphs and graphs with at most two corner chains,
and then reduce the general case to smaller graphs by removing part of the
last clique or the first vertex.

Throughout, $\Gamma$ is an ordinary decorated unit-interval graph, and we only
consider orientations in which every corner is a cover.
In particular every corner is an edge of $\theta\cup S$, and by
Lemma~\ref{lem:decoration}(i) we may assume that all corners are strict.
We write $uv\in\theta$ if $uv$ is an edge of $\theta\cup S$, and covers always
refer to $\Pplus{\theta}$.
We write
\[
 vy\mapsto yh\quad[vh]
\]
for the flip which replaces the edge $vy$ by $yh$ and retains the long edge
$vh$, and similarly when the short edges are exchanged in the other direction.
By Lemma~\ref{lem:flipCriterion}, such an exchange is a flip as soon as the two
edges $vy$ and $vh$ are covers before it, and then $vh$ and $yh$ are covers
after it.
We repeatedly use the following observation: a directed path in $\theta\cup S$
between two vertices of an interval stays in that interval, so the cover test
of an edge $uv$ only involves edges inside $[u,v]$.

\subsection{Corner chains}\label{subs:preparation}

Distinct corners have distinct tails and distinct heads, so the corners form
disjoint directed paths, the \defin{corner chains}; a vertex on no corner is a
chain by itself.
Let $\defin{r(u)}$ be the last vertex of the corner chain containing $u$.
A natural edge $uv$ with $u$ and $v$ on the same chain is a corner: if $uh$ is
the corner leaving $u$, then $v\leq h$ by maximality of $uh$, and $v\geq h$
since $v$ comes after $u$ on the chain.
Let $\Gplus$ be the set of corners together with the normal edges $uv$ with
$r(u)<r(v)$.

If $\theta\subseteq\Gplus$, that is, every edge of $\theta\cup S$ lies in
$\Gplus$, then every corner is a cover.
Indeed, $r$ is constant along corners and increases along the other edges of
$\Gplus$, so a path between the endpoints $u,h$ of a corner only uses corners,
and the only corner leaving $u$ is $uh$.

\begin{lemma}\label{lem:preparation}
Every orientation $\theta$ in which all corners are covers is joined by flips to
an orientation $\theta'\subseteq\Gplus$ which contains every edge of
$\theta\cap\Gplus$.
\end{lemma}

\begin{proof}
We process the corner chains in decreasing order of their last vertices.
Let $U$ be the union of the chains already processed.
We maintain that no edge of $\theta$ leaves $U$, and that all corners are covers
when we start a new chain; initially $U=\emptyset$.

Let $B$ be the next chain, with last vertex $t$.
Consider the edges $vy\in\theta$ with $v\in B$ and $r(y)<t$.
For such an edge, $v<y\leq r(y)<t$, so $v$ has a corner successor $h$ on $B$, and
$y<h$ since $vy$ does not contain the corner $vh$.
Choose a vertex $v$ with such an edge, and among its edges of this kind the edge
$vy$ with the least head.
Then $vy$ is a cover.
Indeed, a path from $v$ to $y$ cannot pass through $U$, since no edge leaves
$U$.
An edge from $v$ to a vertex $z\notin U$ has $r(z)\leq t$, by the choice of $t$.
If $r(z)=t$, it is the corner $vh$, and $h>y$; if $r(z)<t$, it is one of the
edges considered, and
$z\geq y$ by the choice of $vy$.
So the only path from $v$ to $y$ is the edge $vy$.
Since $vh$ is a cover as well, we can perform the flip $vy\mapsto yh\ [vh]$,
\[
 \EdgePic{3}{v,y,h}{1/3/s,1/2/b}\;\longrightarrow\;\EdgePic{3}{v,y,h}{1/3/s,2/3/b}\,,
\]
after which $vh$ is still a cover.
The new edge $yh$ lies inside the corner $vh$, so it is normal, and
$r(y)<r(h)=t$, so it lies in $\Gplus$ and is not one of the edges considered.
The intervals of the other corners of $B$ have interiors disjoint from that of
$vh$, so their cover tests are unchanged, and the edges inside $U$ are
unchanged.
Each flip removes one of the edges considered and creates none, so after
finitely many flips there are none left.

Then no edge leaves $U\cup B$, and every edge created while processing $B$ ends
in $B$.
A path bypassing a corner outside $U\cup B$ cannot use such an edge, since it
could not leave $U\cup B$ again; hence such a bypass would have existed before
we processed $B$.
Thus all corners are covers again.
When all chains have been processed, all edges of the orientation lie in
$\Gplus$, and only edges outside $\Gplus$ were removed.
\end{proof}

\subsection{Exchanging two edges}\label{subs:slides}

The next lemma exchanges two edges along two crossing covers, restoring all
other edges.

\begin{lemma}\label{lem:crossingSlide}
Let $a<b<c<d$, and let $\theta$ be an orientation in which $ac$ and $bd$ are
covers, $ab\in\theta$ and $cd\notin\theta$.
Then $\theta$ is joined by flips to the orientation $\theta-ab+cd$, in which $ac$
and $bd$ are again covers.
\end{lemma}

\begin{proof}
All edges below lie inside $ac$ or $bd$, so they are natural and normal.
We perform the flips of Figure~\ref{fig:crossingSlide} in four stages:
\begin{enumerate}[(1)]
\item $av\mapsto vc\ [ac]$ for all $av\in\theta$ with $a<v<b$, in increasing order of $v$;
\item $wd\mapsto bw\ [bd]$ for all $wd\in\theta$ with $c<w<d$, in decreasing order of $w$;
\item $ab\mapsto bc\ [ac]$ and then $bc\mapsto cd\ [bd]$;
\item the flips of (2) in reverse, and then those of (1) in reverse.
\end{enumerate}
By Lemma~\ref{lem:flipCriterion}, each step is a flip as soon as its two edges
are covers before it, and then the retained cover stays a cover.
In (1), when $av$ is moved the edges from $a$ to smaller vertices have already
been moved, so $av$ is a cover; these flips only change edges with tails below
$b$, so $bd$ stays a cover.
Stage (2) is symmetric, and only changes edges with heads above $c$.
In (3), $ab$ is a cover since $a$ has no edge into $(a,b)$, and $bc\notin\theta$
since otherwise $ab,bc$ would bypass $ac$; after the first flip, $bd$ is still a
cover, since a bypass would begin with $bc$ and no path leads from $c$ to $d$.
In (4), the net change from $ab$ to $cd$ does not affect the edges inside
$[b,w]$ and $[v,c]$, so every edge is a cover when it is moved back, as after
the corresponding flip in (2) or (1).
\end{proof}

\begin{figure}[!ht]
\centering
\begin{gather*}
 \EdgePic[1.55em]{6}{a,v,b,c,w,d}{1/4/c,3/6/c,1/3/b,1/2/b,5/6/b}\xrightarrow{(1)}\EdgePic[1.55em]{6}{a,v,b,c,w,d}{1/4/c,3/6/c,1/3/b,2/4/b,5/6/b}\xrightarrow{(2)}\EdgePic[1.55em]{6}{a,v,b,c,w,d}{1/4/c,3/6/c,1/3/b,2/4/b,3/5/b}\\
 \xrightarrow{(3)}\EdgePic[1.55em]{6}{a,v,b,c,w,d}{1/4/c,3/6/c,3/4/b,2/4/b,3/5/b}\xrightarrow{(3)}\EdgePic[1.55em]{6}{a,v,b,c,w,d}{1/4/c,3/6/c,4/6/b,2/4/b,3/5/b}\xrightarrow{(4)}\EdgePic[1.55em]{6}{a,v,b,c,w,d}{1/4/c,3/6/c,4/6/b,1/2/b,5/6/b}
\end{gather*}
\caption{The flips in the proof of Lemma~\ref{lem:crossingSlide}, with one
edge $av$ with $a<v<b$ and one edge $wd$ with $c<w<d$.
The gray arcs above are the covers $ac$ and $bd$, which are retained
throughout, the arcs below are the edges which move, and the dotted line stands
for all other edges, which are restored at the end.}
\label{fig:crossingSlide}
\end{figure}



We now use Lemma~\ref{lem:crossingSlide} to move edges between two corner chains.
Let $uv$ be a normal natural edge such that $u$ has a corner successor $h$.
Then $v<h$, and $vh$ is natural since it lies inside $uh$.
We call $vh$ the successor of the edge $uv$.
Passing to the successor exchanges the two values of $r$ at the endpoints, and
both endpoints increase.
Every normal edge has at most one successor and at most one predecessor, so the
normal edges form disjoint sequences in which each edge is followed by its
successor.

\begin{lemma}\label{lem:linkPaths}
For two corner chains there is at most one such sequence whose edges join the
two chains.
Along this sequence, the edges in $\Gplus$ and the edges not in $\Gplus$
alternate, and two consecutive edges in $\Gplus$ have the form $ab$ and $cd$
with $a<b<c<d$ such that $ac$ and $bd$ are corners.
\end{lemma}

\begin{proof}
Let $uv$ be the first edge of such a sequence.
If $v$ had a corner predecessor $p$, then $p<u$, since $uv$ does not contain
the corner $pv$, and then $pu$ would be a predecessor of $uv$.
Hence $v$ is the first vertex of its chain, and $u$ lies on a chain starting
earlier.
There is at most one such $u$ on a given chain: if $u<u'<v$ were two
candidates, the natural edge $uv$ would contain the corner leaving $u$.
This proves the first claim.
Passing to the successor exchanges the two values of $r$, so membership in
$\Gplus$ alternates.
If $ab$ is followed by two further edges, they use the corners leaving $a$ and $b$, say
$ac$ and $bd$, and the path continues as $ab\to bc\to cd$ with $a<b<c<d$.
\end{proof}

\begin{lemma}\label{lem:parallel}
Let $\theta,\theta'\subseteq\Gplus$ be orientations which have, for every pair
of corner chains, the same number of normal edges joining them.
Then $\theta$ and $\theta'$ lie in the same flip class.
\end{lemma}

\begin{proof}
List the positions in $\Gplus$ joining a given pair of chains in the order of
their sequence.
By Lemmas~\ref{lem:linkPaths} and~\ref{lem:crossingSlide}, we may replace an occupied
position followed by a vacant one by the reverse, restoring all other edges.
The sum of the indices of the occupied positions increases, so we reach the
configuration occupying the last positions of the list, which only depends on
the number of occupied positions.
We do this for every pair of chains, starting from $\theta$ and from $\theta'$.
\end{proof}

\begin{lemma}\label{lem:twoChains}
If there are at most two corner chains, then all orientations in which every
corner is a cover and which have the same number of ascending edges lie in one
flip class.
\end{lemma}

\begin{proof}
By Lemma~\ref{lem:preparation} we may assume that both orientations lie in
$\Gplus$.
With one chain, the only such orientation consists of the corners.
With two chains, every normal edge of $\Gplus$ joins the two chains, and
Lemma~\ref{lem:parallel} applies.
\end{proof}

\subsection{Complete graphs}\label{subs:complete}

In this subsection we label the vertices $0,1,\dotsc,n$, with $n\geq1$, and
consider the complete natural graph.
Its only corner is $0n$.

\begin{lemma}\label{lem:complete}
For the complete graph on $0,1,\dotsc,n$, the connected orientations in which
$0n$ is a cover and which have the same number of ascending edges lie in one
flip class.
\end{lemma}

\begin{proof}
Let $P_n$ consist of $0n$ and the edges $i(i+1)$ for $1\leq i<n$.
We call the other pairs $ij$, with $1\leq i<i+1<j\leq n$, the chords of $P_n$.
We first join every orientation in question to an orientation $P_n\cup X$, where
$X$ is a set of chords, and then show that the orientations $P_n\cup X$ with the
same number of chords lie in one flip class.

\smallskip\noindent\textbf{Step 1: reaching $P_n\cup X$.}
Here $r(0)=r(n)=n$ and $r(v)=v$ otherwise, so $\Gplus$ consists of $0n$ and all
edges on $1,\dotsc,n$.
By Lemma~\ref{lem:preparation} we may assume that $0n$ is the only edge at $0$.
Since $\theta$ is connected, the set $\eta$ of the remaining edges is connected on
$1,\dotsc,n$.
Whenever a vertex $i$ has two outgoing covers $ij$ and $ik$ in $\eta$, with $j<k$,
we perform the flip $ij\mapsto jk\ [ik]$.
The sum of the tails of the edges increases, so this terminates.
Then every vertex has at most one outgoing cover.
The covers form a connected graph, and a cycle in it would have a least vertex
with two outgoing covers, so they form a tree, with unique sink $n$.
Every edge of $\eta$ joins a vertex to one of its ancestors in this tree.

If $i<j$ are two children of $n$, their only outgoing edges are $in$ and $jn$,
and the four flips
\begin{equation}\label{eq:childMerge}
 in\mapsto0i\ [0n],\qquad jn\mapsto 0j\ [0n],\qquad
 0i\mapsto ij\ [0j],\qquad 0j\mapsto jn\ [0n]
\end{equation}
replace $in$ by $ij$, see Figure~\ref{fig:childMerge}.
For the first two, the vertex losing its edge to $n$ becomes a sink, so the new
edge from $0$ is a cover and $0n$ remains one.
At the third flip, $i$ and $j$ are sinks; afterwards $0$ is adjacent only to
the sink $j$ and to $n$, which gives the last flip.
In the tree, the branch $i\to n$ becomes $i\to j\to n$, and every edge still
joins a vertex to an ancestor.
The vertex $n-1$ is always a child of $n$.
Applying~\eqref{eq:childMerge} with $j=n-1$ until $n-1$ is the only child of
$n$, the edges on $1,\dotsc,n-1$ form a connected graph.
Then we perform $(n-1)n\mapsto 0(n-1)\ [0n]$ and apply the construction
inductively on $0,\dotsc,n-1$.
These flips remain flips in the whole graph, since a path between two vertices
below $n$ cannot pass through $n$.
At the end, $0$ is adjacent only to $n-1$ and $n$, and $n-1$ has no outgoing
edge, so we can perform $0(n-1)\mapsto(n-1)n\ [0n]$.
We have reached an orientation $P_n\cup X$.

\begin{figure}[!ht]
\centering
\begin{gather*}
 \EdgePic{4}{0,i,j,n}{1/4/c,2/4/b,3/4/b}\longrightarrow\EdgePic{4}{0,i,j,n}{1/4/c,1/2/b,3/4/b}\longrightarrow\EdgePic{4}{0,i,j,n}{1/4/c,1/2/b,1/3/b}\\
 \longrightarrow\EdgePic{4}{0,i,j,n}{1/4/c,2/3/b,1/3/b}\longrightarrow\EdgePic{4}{0,i,j,n}{1/4/c,2/3/b,3/4/b}
\end{gather*}
\caption{The four flips~\eqref{eq:childMerge}, which replace $in$ by $ij$; the
cover $0n$ (gray) is retained throughout.
The children $i$ and $j$ of $n$ become the chain $i\to j\to n$.}
\label{fig:childMerge}
\end{figure}

\smallskip\noindent\textbf{Step 2: working in a smaller interval.}
Let $0<p\leq m\leq n$, and consider two orientations $P_n\cup X$ and
$P_n\cup X'$ which agree on all pairs not contained in $W=\{0,p,p+1,\dotsc,m\}$.
We show that a path of flips between their restrictions to $W$, in the complete
graph on $W$, lifts to a path between them.
First we open the chain by the flips
\begin{equation}\label{eq:opening}
 j(j+1)\mapsto 0j\quad[0(j+1)],\qquad j=n-1,n-2,\dotsc,m.
\end{equation}
Before such a flip, all edges from $0$ end at $j+1$ or above, so $0(j+1)$ is a
cover, and these are flips regardless of the chords.
Now the edges at $0$ are $0m,0(m+1),\dotsc,0n$, and the restriction to $W$ is
of the same form for the complete graph on $W$.
A path between two vertices of $W$ stays in $W$: it cannot use a vertex below
$p$, since no edge from $0$ goes there, and it cannot return from above $m$.
Hence flips in the complete graph on $W$ are flips of the whole graph.
Opening $P_n\cup X'$ in the same way gives its restriction to $W$, with the same
edges outside $W$, so after the path inside $W$ we close the chain by reversing
the opening of $P_n\cup X'$.

\smallskip\noindent\textbf{Step 3: rearranging the chords.}
We show by induction on $n$ that the orientations $P_n\cup X$ with the same
number of chords lie in one flip class.
For $n\leq3$ there is at most one chord.
For $n=4$ the chords are $13$, $14$ and $24$.
In the following table, a triple $ijk$ denotes the flip exchanging $ij$ and
$jk$ across $ik$, in the direction allowed by the current state, and each row
starts and ends with $P_4$ together with the indicated chords:
\begin{center}
\begin{tabular}{@{}ll@{}}
\toprule
Chords & Flips \\
\midrule
$\{13\}\to\{14\}$ & $034,023,013,014,023,034$ \\
$\{14\}\to\{24\}$ & $034,124,014,023,012,023,034$ \\
$\{13,24\}\to\{14,24\}$ & $034,024,012,014,024,013,023,012,023,034$ \\
$\{14,24\}\to\{13,14\}$ & $034,023,012,013,023,124,034$ \\
\bottomrule
\end{tabular}
\end{center}
The cover conditions are checked directly at each step.
These rows connect all sets of one or two chords.

Let $n\geq5$.
By Step~2 and induction applied to $W=\{0,1,\dotsc,n-1\}$ and to
$W=\{0,2,3,\dotsc,n\}$, we can rearrange the chords inside either set at a fixed
number, restoring all other edges.
The two sets of chord positions cover all positions except $1n$, and both
contain $24$.
Hence we can exchange any two chords other than $1n$: if no single set contains
both, we route through $24$, moving the selected chord to $24$ first if $24$ is
vacant, and moving $24$ to the vacant position first otherwise.
This connects all orientations $P_n\cup X$ with the same number of chords and the same
status of $1n$.

It remains to change the status of $1n$.
Suppose $1n\notin\theta$ and $13,2n\in\theta$.
We perform $in\mapsto 0i\ [0n]$ for all $in\in\theta$ with $0<i<n$, in decreasing
order of $i$.
When $in$ is moved, the edges into $n$ with larger tails have been moved.
Hence $in$ is a cover, and the vertices other than $n$ adjacent to $0$, which are
larger than $i$, cannot reach $n$, so $0n$ is a cover.
Now $02\in\theta$ and $01\notin\theta$, so $02$ is a cover, and we perform the
flip $12\mapsto 01\ [02]$.
Then we reverse these flips in increasing order of $i$, including the new edge
$01$.
When $0i$ is moved back to $in$, it has the least head among the edges at $0$,
and the edges into $n$ restored so far have smaller tails, so $0i$ and $0n$ are
covers.
We obtain $\theta-12+1n$.
To restore $12$, we open the chain down to $3$ by~\eqref{eq:opening}.
On $0,1,2,3$ the edges are then $03$, $13$ and $23$, and~\eqref{eq:childMerge}
with $(i,j,n)=(1,2,3)$ replaces $13$ by $12$.
Closing the chain gives
\begin{equation}\label{eq:outerChord}
 \theta\longmapsto\theta-13+1n\qquad(13,2n\in\theta,\ 1n\notin\theta).
\end{equation}
Let $\theta=P_n\cup X$; there are $\binom{n-1}{2}$ chords in total.
If $2\leq|X|\leq\binom{n-1}{2}-1$, we can first arrange $13,2n\in\theta$ and
$1n\notin\theta$, and~\eqref{eq:outerChord} connects the two statuses of $1n$.
If $|X|=0$ or $|X|=\binom{n-1}{2}$, there is only one set of chords.
If $|X|=1$, start from $P_n+13$, open the chain down to $2$, perform
$13\mapsto01\ [03]$ and $01\mapsto 1n\ [0n]$, and close the chain.
While the chain is open, the vertices $2,\dotsc,n$ are sinks, so these are flips,
and we reach $P_n+1n$.
This completes the induction, and the proof.
\end{proof}

\subsection{Deleting a chain or an edge}

In this subsection, connectivity means connectivity of undirected graphs.
Suppose that the natural graph is connected, has at least three vertices, and
has no articulation.
Then every pair $uv$ with $v-u\leq2$ is a natural edge.
Indeed, if $u(u+1)$ were not natural, then by closure under subintervals no
natural edge would join $[1,u]$ and $[u+1,n]$; and if $u(u+2)$ were not natural,
no natural edge would pass over $u+1$, which would be an articulation.

Replacing the natural graph by $\Gplus$ preserves connectivity within any union
of corner chains.
Indeed, if $uv$ is normal with $r(u)>r(v)$, and $uh$ is the corner leaving $u$,
then $u<v<h$, the edge $vh$ lies in $\Gplus$, and the walk $u,h,v$ only visits
the chains of $u$ and $v$.

\begin{lemma}\label{lem:chainDeletion}
Deleting the vertices of one corner chain leaves the other vertices connected
in $\Gplus$.
\end{lemma}

\begin{proof}
Two consecutive vertices $v,v+1$ are not on the same chain, since their edge
would be a corner, while it lies inside $(v-1)(v+1)$ or $v(v+2)$.
Hence the remaining vertices, listed in increasing order, have differences at
most two, and consecutive ones are joined by natural edges.
We replace the edges outside $\Gplus$ by the walks above, which avoid the deleted
chain.
\end{proof}

\begin{lemma}\label{lem:removableEdge}
If there are at least three corner chains, then deleting any normal edge $uv$ of
$\Gplus$ leaves $\Gplus$ connected.
\end{lemma}

\begin{proof}
Choose a vertex $w$ on a third chain.
By Lemma~\ref{lem:chainDeletion}, there is a path from $u$ to $w$ avoiding the
chain of $v$, and a path from $w$ to $v$ avoiding the chain of $u$.
Neither uses $uv$.
\end{proof}

\subsection{The last clique}

Suppose that the last two corners are $bc$ and $an$, with $b<a<c<n$.
Let $L=[1,c]\cup\{n\}$, and let $\Gamma[L]$ be the graph induced on $L$, with the
inherited order and strict edges.
For an orientation $\theta$ we write $\theta[L]$ for its edges inside $L$.
All corners other than $an$ end at most at $c$, so the vertices in $(c,n)$ lie
on no corner.
Hence $\Gamma[L]$ has the same corners as $\Gamma$, its corner chains are those of
$\Gamma$ except for the singleton chains of the vertices in $(c,n)$, and its
set $\Gplus$ consists of the edges of $\Gplus$ inside $L$.
Every vertex strictly between $a$ and $n$ is the last vertex of its chain.
Hence an orientation in $\Gplus$ has no edge from $a$ other than $an$, and
every natural edge with both endpoints strictly between $a$ and $n$, or ending
at $n$ and starting after $a$, lies in $\Gplus$.

\begin{lemma}\label{lem:lastClique}
If Theorem~\ref{thm:connectivity} holds for $\Gamma[L]$, then it holds for
$\Gamma$.
\end{lemma}

\begin{proof}
Let $\theta$ be a connected orientation in which all corners are covers, and
assume $\theta\subseteq\Gplus$ by Lemma~\ref{lem:preparation}.
We join $\theta$ to an orientation containing a fixed spanning tree $T$, and then
join all such orientations of the same size; Figure~\ref{fig:lastCliqueFlips}
shows the flips used.

All pairs in $[a,n]$ are natural, and a path between two vertices of a connected
component of the edges of $\theta$ inside $[a,n]$ stays in that component.
So flips of the complete graph on such a component are flips of $\Gamma$, and
Lemma~\ref{lem:complete} applies to the component containing $a$ and $n$, since
$an$ is a cover.

\smallskip\noindent\textbf{Step 1: making $\theta[L]$ connected.}
In the component containing $a$ and $n$, the vertex $a$ is only adjacent to $n$,
and by Lemma~\ref{lem:complete} we replace its edges by the star of all edges to
$n$, together with edges avoiding $a$; this is possible since a connected graph
on $k$ vertices with a leaf has between $k-1$ and $1+\binom{k-1}{2}$ edges.
Every other component meets $(a,c]$, since otherwise it would be a connected
component of $\theta$.
Join its least vertex $u$ to another vertex $v\leq c$ by a path of covers, using
edges in either direction.
If the largest vertex $z$ of the path exceeds $v$, its neighbours $i<j$ give the
flip $jz\mapsto ij\ [iz]$ of Figure~\ref{fig:lastCliqueFlips}(a), which shortens the
path; replacing an edge which is no longer a cover by a path of covers inside its
interval, the largest vertex decreases.
We end with a path inside $[u,v]$, and these flips keep all pairs joined inside
$[1,c]$.
Every natural edge with an endpoint in $(c,n)$ lies inside $an$, so $\theta[L]$
is now connected.

\smallskip\noindent\textbf{Step 2: changing $\theta[L]$.}
The flips $vn\mapsto av\ [an]$ of Figure~\ref{fig:lastCliqueFlips}(b), for
$vn\in\theta$ with $c<v<n$ in decreasing order of $v$, are flips as in stage~(1)
of the proof of Lemma~\ref{lem:crossingSlide}, and they change neither the edges
inside $L$ nor the cover tests of other corners.
Afterwards a path leaving $L$ cannot return to $L$, so every path of flips in
$\Gamma[L]$ lifts, with all edges outside $L$ fixed.
By assumption we reach any connected orientation of $\Gamma[L]$ in $\Gplus$ with
the same number of edges, and then we undo the flips of
Figure~\ref{fig:lastCliqueFlips}(b) in increasing order of $v$.

\smallskip\noindent\textbf{Step 3: a fixed tree.}
In the set $\Gplus$ of $\Gamma[L]$, the corners and the edges $vn$ with $a<v\leq c$
form a forest, since every vertex has at most one outgoing edge among them.
It extends to a spanning tree $T'$ of $\Gamma[L]$ inside $\Gplus$.
Inside $[a,n]$, the tree $T'$ consists of edges to $n$, and $ba\notin T'$ since
$bc,cn,na$ already join $b$ and $a$.
By Step~2 we may assume that $T'\subseteq\theta$.
Then $[a,n]$ is one component: the vertices in $(c,n)$ are joined to $L$ by
unchanged edges, and those of $L\cap[a,n]$ are joined to $n$ by $T'$.
Lemma~\ref{lem:complete} makes it the star of Figure~\ref{fig:lastCliqueFlips}(c)
together with edges avoiding $a$, and we obtain an orientation containing
$T=T'\cup\{vn:c<v<n\}$, see Figure~\ref{fig:lastCliqueTree}.

\begin{figure}[!ht]
\centering
\begin{gather*}
 \text{(a)}\quad \EdgePic[2em]{3}{i,j,z}{1/3/c,2/3/b}\longrightarrow\EdgePic[2em]{3}{i,j,z}{1/3/c,1/2/b}
 \qquad\qquad
 \text{(c)}\quad \EdgePic[1.6em]{5}{a,c,v,w,n}{1/5/c,2/5/b,3/5/b,4/5/b}\\[1ex]
 \text{(b)}\quad \EdgePic[2em]{4}{a,c,v,n}{1/4/c,3/4/b}\longrightarrow\EdgePic[2em]{4}{a,c,v,n}{1/4/c,1/3/b}\\[1ex]
 \text{(d)}\quad \EdgePic[2em]{4}{b,a,c,n}{1/3/c,2/4/c,1/2/b}\longrightarrow\EdgePic[2em]{4}{b,a,c,n}{1/3/c,2/4/c,3/4/b}
\end{gather*}
\caption{The flips in the proof of Lemma~\ref{lem:lastClique}, with retained
covers in gray above the line and moving edges below.
(a) Shortening a path of covers through its largest vertex $z$ in Step~1.
(b) Moving an edge $vn$ with $c<v<n$ off $n$ in Step~2.
(c) The star at $n$ inside $[a,n]$ in Step~3.
(d) Exchanging $ba$ for $cn$ in Step~4, by Lemma~\ref{lem:crossingSlide} for the
corners $bc$ and $an$.}
\label{fig:lastCliqueFlips}
\end{figure}

\begin{figure}[!ht]
\centering
\LLTpic[cell=0.9em]{nnnnddenndeeeee}{\LLTasc{3/5,5/9,6/9}\LLTmark{2/4}{$\circ$}}\qquad\qquad
\LLTpic[cell=0.9em]{nnnnddenndeeeee}{\LLTasc{3/5,5/9,6/9,7/9,8/9}\LLTmark{2/4}{$\circ$}}
\caption{The trees $T'$ (left) and $T$ (right) for the path
$\mathtt{nnnnddenndeeeee}$, with corners $15$, $26$ and $49$, so that $b=2$, $a=4$,
$c=6$, $n=9$ and $L=[1,6]\cup\{9\}$.
The diagonal steps cross the corners, which are tree edges, and the arrows
are the remaining tree edges; $T$ adds the leaves $7$ and $8$ at $9$.
The position $ba=24$, marked $\circ$, lies in neither tree.}
\label{fig:lastCliqueTree}
\end{figure}

\smallskip\noindent\textbf{Step 4: moving the remaining edges.}
Let $M_1$ be the positions of $\Gplus\setminus T$ inside $L$, and let $M_2$ be the
set of pairs $uv$ with $a<u<v<n$, together with $ba$.
Together they are all positions of $\Gplus\setminus T$, and $ba\in M_1\cap M_2$.
Step~2 exchanges any two positions of $M_1$, and Lemma~\ref{lem:complete} on the
component $[a,n]$ exchanges any two positions of $M_2\setminus\{ba\}$, restoring all
other edges.
Suppose $ba\in\theta$ and $x\in M_2\setminus\{ba\}$ is vacant.
Let $d\neq c$ be an endpoint of $x$, and let $y$ be the edge joining $c$ and $d$.
If $y\in\theta$, the component stays connected when we replace $cn$ by $x$, so
Lemma~\ref{lem:complete} performs this replacement.
Then Lemma~\ref{lem:crossingSlide} for the corners $bc$ and $an$ exchanges $ba$
for $cn$, see Figure~\ref{fig:lastCliqueFlips}(d),
and we obtain $\theta-ba+x$.
If $y\notin\theta$, the same construction with $y$ in place of $x$ gives
$\theta-ba+y$, and then we exchange $y$ and $x$.
Reversing these paths exchanges in the other direction.
Finally, a position of $M_1\setminus M_2$ and one of $M_2\setminus M_1$ are
exchanged by routing through $ba$, as in Step~3 of Lemma~\ref{lem:complete}.
Thus all orientations containing $T$ with the same number of edges lie in one flip
class, and with Lemma~\ref{lem:preparation} this proves the lemma.
\end{proof}

Reversing the order of the vertices maps flips to flips, since it reverses the
direction of all edges and preserves covers.
Hence Lemma~\ref{lem:lastClique} also holds for the first two corners, with the
graph induced on $\{1\}\cup[b',n]$ in place of $\Gamma[L]$ when the first two
corners are $1a'$ and $b'c'$.

\subsection{The first vertex}

Suppose that the first two corners are $1a$ and $2b$, with $2<a<b$.
Let $\Gamma'$ be the graph induced on $2,\dotsc,n$.
It has the corners of $\Gamma$ other than $1a$, and the same values of $r$ on
$2,\dotsc,n$.
Hence $\Gplus$ consists of $1a$, the set $(\Gamma')^{+}$ for $\Gamma'$, and the
edges $1v$ for $v$ in
\[
 J=\{v:1<v<a,\ r(a)<r(v)\}.
\]
For $v\in J$, let $vh_v$ be the corner leaving $v$.
Then $h_v>a$, since a corner starting in $(1,a)$ does not end inside $[1,a]$,
and $ah_v$ is a normal edge of $(\Gamma')^{+}$.
Lemma~\ref{lem:crossingSlide} for the corners $1a$ and $vh_v$ exchanges $1v$ and
$ah_v$:
\[
 \EdgePic{4}{1,v,a,h_v}{1/3/s,2/4/s,1/2/b}
 \;\longleftrightarrow\;
 \EdgePic{4}{1,v,a,h_v}{1/3/s,2/4/s,3/4/b}\,.
\]

\begin{lemma}\label{lem:firstVertex}
Suppose that Theorem~\ref{thm:connectivity} holds for $\Gamma'$, and that
$(\Gamma')^{+}\setminus\{ah_v\}$ is connected for every $v\in J$.
Then Theorem~\ref{thm:connectivity} holds for $\Gamma$.
\end{lemma}

\begin{proof}
Let $\theta\subseteq\Gplus$ be connected.
If the edges of $\theta$ inside $[2,n]$ do not form a connected graph, let $B$ be the
component containing $a$.
Some edge $1v\in\theta$ ends in another component, and then $v\in J$, the vertices
$v$ and $h_v$ lie in the same component, and $ah_v\notin\theta$.
Lemma~\ref{lem:crossingSlide} replaces $1v$ by $ah_v$, and the number of edges at $1$
decreases.
Hence we may assume that the edges inside $[2,n]$ form a connected graph.

Fix the set $I\subseteq J$ of vertices $v$ with $1v\in\theta$.
Paths between vertices of $[2,n]$ avoid $1$, so flips of $\Gamma'$ lift with the
edges at $1$ fixed, and by assumption all such orientations with the same $I$
and the same number of edges lie in one flip class.

Let $m$ be the number of edges of $(\Gamma')^{+}$, and let $s$ be the number of
edges of $\theta$.
The corners of $\Gamma'$ form a forest, which extends to a spanning tree of
$(\Gamma')^{+}$, so connected orientations of $\Gamma'$ in $(\Gamma')^{+}$ exist
with any number of edges between $n-2$ and $m$.
Hence the possible sets are exactly the $I\subseteq J$ with
\begin{equation}\label{eq:initialBand}
 \max(0,s-1-m)\leq|I|\leq\min(|J|,s-1-(n-2)).
\end{equation}
Suppose that $I$ and $I\setminus\{v\}$ both satisfy~\eqref{eq:initialBand}.
The corresponding orientations of $\Gamma'$ have $t$ and $t+1$ edges, with
$n-2\leq t<m$.
Since $(\Gamma')^{+}\setminus\{ah_v\}$ is connected, there is a connected
orientation of $\Gamma'$ in $(\Gamma')^{+}$ with $t$ edges avoiding $ah_v$.
Together with the edges $1w$, $w\in I$, Lemma~\ref{lem:crossingSlide} moves it
to an
orientation with the set $I\setminus\{v\}$.
Thus the classes for $I$ and $I\setminus\{v\}$ agree.

Any two sets satisfying~\eqref{eq:initialBand} are joined by such insertions and
deletions: to exchange two elements we delete first if $|I|$ is above its lower
bound, and insert first otherwise.
This fails only if~\eqref{eq:initialBand} allows a single cardinality and more
than one set.
If $J\neq\emptyset$, the hypothesis gives $m\geq n-1$, and then a single allowed
cardinality forces $I=\emptyset$ or $I=J$.
Hence all connected orientations in $\Gplus$ with $s$ edges lie in one flip
class, and Lemma~\ref{lem:preparation} proves the lemma.
\end{proof}

\subsection{Proof of Theorem~\ref{thm:connectivity}}

First let the natural graph be connected without articulations.
We proceed by induction on $n$.
If there is only one corner, the graph is complete, and Lemma~\ref{lem:complete}
applies.
If there are at most two corner chains, Lemma~\ref{lem:twoChains} applies.
Otherwise let $1a$ and $b'c'$ be the first two corners.
Since there are no articulations, consecutive maximal cliques share at least two
vertices, so $1<b'<a<c'$.
If $b'>2$, the reversed form of Lemma~\ref{lem:lastClique} reduces the claim to the
graph induced on $\{1\}\cup[b',n]$, which has fewer vertices and no
articulation.
If $b'=2$, deleting the vertex $1$ removes the corner $1a$ and gives a graph
$\Gamma'$ without articulations, whose maximal cliques start with $[2,c']$.
It has the same number of corner chains as $\Gamma$, at least three, so
Lemma~\ref{lem:removableEdge} for $\Gamma'$ verifies the hypothesis of
Lemma~\ref{lem:firstVertex}, and induction applies to $\Gamma'$.

Finally, let the natural graph consist of the path $1,\dotsc,b$ followed by a
graph without articulations on $[b,n]$.
The path edges are corners, so they lie in every orientation we consider, and
paths between vertices of $[b,n]$ do not use vertices below $b$.
Hence restriction to $[b,n]$ is a bijection between the orientations in question
of $\Gamma$ and of the graph induced on $[b,n]$, compatible with flips.
This completes the proof of Theorem~\ref{thm:connectivity}, and thereby of
Theorem~\ref{thm:main} for ordinary graphs.
